\documentclass[12pt,a4paper]{article}

\usepackage[utf8]{inputenc}
\usepackage[T1]{fontenc}
\usepackage[ngerman]{babel}
\usepackage{geometry}
\usepackage{xcolor}
\usepackage{hyperref}
\usepackage{longtable}
\usepackage{array}
\usepackage{listings}
\usepackage{enumitem}
\usepackage{titlesec}
\usepackage{fancyhdr}

\geometry{left=2.5cm,right=2.5cm,top=2.5cm,bottom=2.5cm}
\hypersetup{colorlinks=true,linkcolor=blue,urlcolor=blue}

\definecolor{codebg}{RGB}{245,247,250}
\definecolor{codeframe}{RGB}{210,220,230}
\definecolor{darkblue}{RGB}{20,55,100}

\lstdefinestyle{codestyle}{
    backgroundcolor=\color{codebg},
    frame=single,
    rulecolor=\color{codeframe},
    basicstyle=\ttfamily\small,
    breaklines=true,
    tabsize=4,
    showstringspaces=false
}
\lstset{style=codestyle}

\titleformat{\section}{\Large\bfseries\color{darkblue}}{\thesection}{1em}{}
\titleformat{\subsection}{\large\bfseries\color{darkblue}}{\thesubsection}{1em}{}

\pagestyle{fancy}
\fancyhf{}
\lhead{Crystal Caverns}
\rhead{Projektdokumentation}
\cfoot{\thepage}

\title{\textbf{Crystal Caverns}\\\large Projektdokumentation}
\author{C++17 / SDL2 Spielprojekt}
\date{\today}

\begin{document}
\maketitle
\thispagestyle{empty}

\begin{center}
\vspace{1cm}
\textbf{Repository:} \\ \url{https://github.com/Mc-smail/CrystalCaverns}
\end{center}

\newpage
\tableofcontents
\newpage

\section{Kurzbeschreibung}

\textbf{Crystal Caverns} ist ein kleines 2D-Retro-Spiel in C++17 mit SDL2. Der Spieler bewegt sich durch eine Höhle, sammelt Kristalle, vermeidet Gegner und fallende Steine und erreicht den Ausgang. Nach Abschluss des ersten Levels startet automatisch ein schwierigeres zweites Level.

Das Spiel demonstriert zentrale Konzepte der 2D-Spieleentwicklung:

\begin{itemize}
    \item Game Loop: Input, Update und Render
    \item SDL2-Fenster und Renderer
    \item Event-Handling und Tastatursteuerung
    \item Tilemap aus Textdateien
    \item Sprites und Texturen mit SDL
    \item Kollisionen
    \item objektorientierte Programmierung mit Klassen und Vererbung
    \item moderne Speicherverwaltung mit \texttt{std::unique\_ptr}
\end{itemize}

\section{Spielziel}

Der Spieler muss pro Level mindestens \textbf{7 Kristalle} sammeln. Danach öffnet sich der Ausgang. Wenn der Spieler den Ausgang erreicht, wird im ersten Level automatisch das zweite Level geladen. Nach Abschluss des zweiten Levels ist das Spiel gewonnen.

Der Spieler verliert ein Leben, wenn:

\begin{itemize}
    \item er einen Gegner berührt,
    \item er von einem fallenden Stein getroffen wird,
    \item der Level-Timer abläuft.
\end{itemize}

Der Spieler besitzt insgesamt \textbf{3 Leben}. Bei 0 Leben endet das Spiel mit Game Over.

\section{Steuerung}

\begin{longtable}{|p{5cm}|p{8cm}|}
\hline
\textbf{Taste} & \textbf{Funktion} \\
\hline
\texttt{W} oder Pfeil hoch & Spieler nach oben bewegen \\
\hline
\texttt{A} oder Pfeil links & Spieler nach links bewegen \\
\hline
\texttt{S} oder Pfeil runter & Spieler nach unten bewegen \\
\hline
\texttt{D} oder Pfeil rechts & Spieler nach rechts bewegen \\
\hline
\texttt{R} & Neustart nach Game Over oder Victory \\
\hline
\texttt{ESC} & Spiel beenden \\
\hline
\end{longtable}

\section{Spielmechaniken}

\subsection{Tilemap}

Die Spielwelt besteht aus einem Raster. Jedes Zeichen in der Leveldatei steht für ein bestimmtes Objekt.

\begin{longtable}{|p{3cm}|p{10cm}|}
\hline
\textbf{Zeichen} & \textbf{Bedeutung} \\
\hline
\texttt{\#} & Wand, nicht begehbar \\
\hline
\texttt{.} & Erde bzw. begehbares Feld \\
\hline
\texttt{P} & Startposition des Spielers \\
\hline
\texttt{C} & Kristall \\
\hline
\texttt{O} & Stein bzw. Boulder \\
\hline
\texttt{X} & Gegner \\
\hline
\texttt{E} & Ausgang \\
\hline
\end{longtable}

\subsection{Kristalle und Ausgang}

Kristalle werden eingesammelt, wenn der Spieler das entsprechende Feld betritt. Sobald mindestens 7 Kristalle gesammelt wurden, öffnet sich der Ausgang. Erreicht der Spieler den geöffneten Ausgang, wird entweder das nächste Level geladen oder das Spiel gewonnen.

\subsection{Gegner}

Gegner bewegen sich automatisch horizontal. Wenn ein Gegner auf eine Wand oder ein Hindernis trifft, wechselt er seine Richtung. Berührt der Spieler einen Gegner, verliert er ein Leben.

\subsection{Steine}

Steine fallen nach unten, wenn das Feld darunter frei ist. Wird der Spieler von einem fallenden Stein getroffen, verliert er ein Leben.

\subsection{Timer und Leben}

Jedes Level hat ein Zeitlimit von \textbf{90 Sekunden}. Das HUD zeigt Kristalle, Leben und eine Timer-Leiste an.

\section{Levelübersicht}

\begin{longtable}{|p{3cm}|p{3cm}|p{3cm}|p{3cm}|p{3cm}|}
\hline
\textbf{Level} & \textbf{Kristalle} & \textbf{Gegner} & \textbf{Steine} & \textbf{Schwierigkeit} \\
\hline
Level 1 & 8 & 3 & 4 & Mittel \\
\hline
Level 2 & 10 & 3 & 5 & Schwer \\
\hline
\end{longtable}

Level 2 ist schwieriger, weil es mehr Kristalle, mehr Steine und engere Wege enthält.

\section{Technische Architektur}

Das Projekt ist objektorientiert aufgebaut. Die wichtigsten Klassen sind:

\begin{longtable}{|p{4cm}|p{10cm}|}
\hline
\textbf{Klasse} & \textbf{Aufgabe} \\
\hline
\texttt{Game} & Hauptklasse, Game Loop, Levelwechsel, Sieg und Niederlage \\
\hline
\texttt{InputManager} & Tastatur- und SDL-Event-Verarbeitung \\
\hline
\texttt{TextureManager} & Laden und Verwalten von BMP-Texturen \\
\hline
\texttt{Tile} & Einzelnes Feld der Map \\
\hline
\texttt{TileMap} & Lädt und rendert das Level \\
\hline
\texttt{Entity} & Abstrakte Basisklasse für bewegliche Objekte \\
\hline
\texttt{Player} & Spielerbewegung und Kristalle sammeln \\
\hline
\texttt{Enemy} & Automatisch bewegter Gegner \\
\hline
\texttt{Boulder} & Fallender Stein \\
\hline
\end{longtable}

\subsection{Game Loop}

Die Hauptschleife besteht aus drei Phasen:

\begin{enumerate}
    \item \textbf{Input}: Tastatur und SDL-Events lesen.
    \item \textbf{Update}: Weltzustand, Gegner, Steine, Timer und Kollisionen aktualisieren.
    \item \textbf{Render}: Tilemap, Spieler, Gegner, Steine und HUD zeichnen.
\end{enumerate}

\begin{lstlisting}[language=C++]
while (running) {
    handleEvents();
    update(deltaTime);
    render();
}
\end{lstlisting}

\section{UML-Klassendiagramm}

\begin{lstlisting}
Game --> InputManager
Game --> TextureManager
Game --> TileMap
Game --> Player
Game --> Enemy
Game --> Boulder
TileMap --> Tile
Entity <|-- Player
Entity <|-- Enemy
Entity <|-- Boulder
\end{lstlisting}

\section{Projektstruktur}

\begin{lstlisting}
CrystalCaverns/
├── CMakeLists.txt
├── README.md
├── assets/
│   ├── levels/
│   │   ├── level01.txt
│   │   └── level02.txt
│   └── sprites/
│       ├── boulder.bmp
│       ├── crystal.bmp
│       ├── dirt.bmp
│       ├── empty.bmp
│       ├── enemy.bmp
│       ├── exit_closed.bmp
│       ├── exit_open.bmp
│       ├── player.bmp
│       └── wall.bmp
├── include/
│   ├── Boulder.hpp
│   ├── Enemy.hpp
│   ├── Entity.hpp
│   ├── Game.hpp
│   ├── InputManager.hpp
│   ├── Player.hpp
│   ├── TextureManager.hpp
│   ├── Tile.hpp
│   ├── TileMap.hpp
│   └── Types.hpp
└── src/
    ├── Boulder.cpp
    ├── Enemy.cpp
    ├── Entity.cpp
    ├── Game.cpp
    ├── InputManager.cpp
    ├── main.cpp
    ├── Player.cpp
    ├── TextureManager.cpp
    ├── Tile.cpp
    └── TileMap.cpp
\end{lstlisting}

\section{Build- und Startanleitung}

\subsection{macOS}

\begin{lstlisting}[language=bash]
brew install cmake sdl2

git clone https://github.com/Mc-smail/CrystalCaverns.git
cd CrystalCaverns
mkdir -p build
cd build
cmake ..
cmake --build .
cd ..
./build/CrystalCaverns
\end{lstlisting}

\subsection{Ubuntu, Kali oder Debian}

\begin{lstlisting}[language=bash]
sudo apt update
sudo apt install -y build-essential cmake libsdl2-dev

git clone https://github.com/Mc-smail/CrystalCaverns.git
cd CrystalCaverns
mkdir -p build
cd build
cmake ..
cmake --build .
cd ..
./build/CrystalCaverns
\end{lstlisting}

\section{Grafik und SDL-Rendering}

Das Spiel benutzt echte SDL-Grafikdateien im BMP-Format. Dadurch wird keine zusätzliche Bibliothek wie SDL2\_image benötigt. SDL2 kann BMP-Dateien direkt laden.

\begin{lstlisting}[language=C++]
SDL_Surface* surface = SDL_LoadBMP(path.c_str());
SDL_Texture* texture = SDL_CreateTextureFromSurface(renderer, surface);
SDL_RenderCopy(renderer, texture, nullptr, &rect);
\end{lstlisting}

\section{Designentscheidungen}

\subsection{Warum SDL2?}

SDL2 stellt die grundlegenden Funktionen bereit, die für ein 2D-Spiel nötig sind: Fenster, Renderer, Events, Tastatur und Texturen. Gleichzeitig bleibt SDL2 übersichtlich genug, um die Architektur selbst zu verstehen.

\subsection{Warum BMP-Dateien?}

BMP-Dateien können direkt mit SDL2 geladen werden. Dadurch bleibt das Projekt einfacher, weil keine zusätzliche Grafikbibliothek installiert werden muss.

\subsection{Warum \texttt{std::unique\_ptr}?}

Dynamische Spielobjekte wie Gegner und Steine werden über Smart Pointer verwaltet. Dadurch wird Speicher automatisch freigegeben und das Projekt bleibt sicherer.

\section{Erweiterungsmöglichkeiten}

\begin{itemize}
    \item mehr Level
    \item Menü mit Text über SDL\_ttf
    \item Soundeffekte mit SDL\_mixer
    \item bessere Gegner-KI
    \item Highscore-System
    \item Level-Editor
    \item Animationen mit Spritesheets
    \item Schlüssel-Tür-System
    \item verschiedene Steinarten
    \item zufällig generierte Höhlen
\end{itemize}

\section{Fazit}

Crystal Caverns ist ein kleines, aber vollständiges C++/SDL2-Spielprojekt. Es zeigt Game Loop, Input, Rendering, Tilemaps, Kollisionen, Entities und objektorientierte Struktur in einem praktischen Spiel. Durch zwei Level, Timer, Leben und steigende Schwierigkeit ist das Spiel nicht nur eine technische Demo, sondern tatsächlich spielbar.

\end{document}
